\documentclass[10pt,a4paper]{amsart}
\usepackage{amsmath,amssymb,amsthm,graphicx}
% Retained minimal preamble for the Stacks Project chapter more-algebra.tex (More on Algebra).
% The chapter file pinned-original-source.tex carries only \input{preamble} on its first line, so no
% preamble text lives inside the chapter itself.  The Stacks Project uses one single shared preamble for
% every chapter; the snapshot stages it identically as sources/stacks-chow/preamble.tex and
% sources/stacks-etale-cohomology/preamble.tex.  The block below is copied mechanically, line for line,
% from sources/stacks-chow/preamble.tex lines 14-18 (the Xy-pic setup) and 154-217 (the native theorem
% declarations together with the mathematical macros).  Nothing was added, removed or redefined; the
% wrapper class and the amsmath/amssymb/amsthm/graphicx base packages are supplied by the assembler.
% Chapter: more-algebra.tex -- More on Algebra.
% For commutative diagrams we use Xy-pic
\usepackage[all]{xy}

% We use 2cell for 2-commutative diagrams.
\xyoption{2cell}

% Theorem environments.
%
\theoremstyle{plain}
\newtheorem{theorem}[subsection]{Theorem}
\newtheorem{proposition}[subsection]{Proposition}
\newtheorem{lemma}[subsection]{Lemma}

\theoremstyle{definition}
\newtheorem{definition}[subsection]{Definition}
\newtheorem{example}[subsection]{Example}
\newtheorem{exercise}[subsection]{Exercise}
\newtheorem{situation}[subsection]{Situation}

\theoremstyle{remark}
\newtheorem{remark}[subsection]{Remark}
\newtheorem{remarks}[subsection]{Remarks}

\numberwithin{equation}{subsection}

% Macros
%
\def\lim{\mathop{\mathrm{lim}}\nolimits}
\def\colim{\mathop{\mathrm{colim}}\nolimits}
\def\Spec{\mathop{\mathrm{Spec}}}
\def\Hom{\mathop{\mathrm{Hom}}\nolimits}
\def\Ext{\mathop{\mathrm{Ext}}\nolimits}
\def\SheafHom{\mathop{\mathcal{H}\!\mathit{om}}\nolimits}
\def\SheafExt{\mathop{\mathcal{E}\!\mathit{xt}}\nolimits}
\def\Sch{\mathit{Sch}}
\def\Mor{\mathop{\mathrm{Mor}}\nolimits}
\def\Ob{\mathop{\mathrm{Ob}}\nolimits}
\def\Sh{\mathop{\mathit{Sh}}\nolimits}
\def\NL{\mathop{N\!L}\nolimits}
\def\CH{\mathop{\mathrm{CH}}\nolimits}
\def\proetale{{pro\text{-}\acute{e}tale}}
\def\etale{{\acute{e}tale}}
\def\QCoh{\mathit{QCoh}}
\def\Ker{\mathop{\mathrm{Ker}}}
\def\Im{\mathop{\mathrm{Im}}}
\def\Coker{\mathop{\mathrm{Coker}}}
\def\Coim{\mathop{\mathrm{Coim}}}

% Boxtimes
%
\DeclareMathSymbol{\boxtimes}{\mathbin}{AMSa}{"02}

%
% Macros for moduli stacks/spaces
%
\def\QCohstack{\mathcal{QC}\!\mathit{oh}}
\def\Cohstack{\mathcal{C}\!\mathit{oh}}
\def\Spacesstack{\mathcal{S}\!\mathit{paces}}
\def\Quotfunctor{\mathrm{Quot}}
\def\Hilbfunctor{\mathrm{Hilb}}
\def\Curvesstack{\mathcal{C}\!\mathit{urves}}
\def\Polarizedstack{\mathcal{P}\!\mathit{olarized}}
\def\Complexesstack{\mathcal{C}\!\mathit{omplexes}}
% \Pic is the operator that assigns to X its picard group, usage \Pic(X)
% \Picardstack_{X/B} denotes the Picard stack of X over B
% \Picardfunctor_{X/B} denotes the Picard functor of X over B
\def\Pic{\mathop{\mathrm{Pic}}\nolimits}
\def\Picardstack{\mathcal{P}\!\mathit{ic}}
\def\Picardfunctor{\mathrm{Pic}}
\def\Deformationcategory{\mathcal{D}\!\mathit{ef}}
\title{Building a finite ring extension that realizes a prescribed isomorphism modulo a fixed power of a nonzerodivisor}
\author{The Stacks Project Authors}
\date{stacks-project:more-algebra.tex:local-snapshot-20261012; GNU FDL 1.2 or later}
\begin{document}
\maketitle
\setcounter{subsection}{1}
\noindent\textit{Editorial note.} All delivered mathematical text is the author's own, copied from the pinned Stacks Project chapter file more-algebra.tex; no statement or proof step was rewritten, reordered or removed. The counted statement is the labeled lemma environment at source lines 35690-35700 and the counted proof is the proof environment at source lines 35702-35725, which immediately follows it and is its own proof.\par\medskip
\noindent\textit{Editorial note.} The last sentence of the proof, at source lines 35721-35724, leaves the final compatibility verification to the reader: it asks the reader to see that the displayed isomorphism of modules sends fA into gB and induces a ring isomorphism phi : A/fA -> B/gB compatible with phi', and it adds that the same display even induces an isomorphism A/f\textasciicircum{}tA -> B/g\textasciicircum{}tB which is not needed. That sentence is preserved verbatim inside the counted span and is the only step the author does not spell out; the construction, the ring structure, the displayed isomorphism and the finiteness argument are all carried out in the span.\par\medskip
\noindent\textit{Editorial note.} The author's text contains the typographical slip Cchoose where choose is meant, at source line 35704. It is copied exactly as printed and is not corrected, since the lane delivers the author's own text.\par\medskip
\noindent\textit{Editorial note.} The chapter file pinned-original-source.tex carries only an input of the shared preamble on its first line, so no preamble text lives inside the chapter. The retained file sources/stacks-more-algebra/retained-preamble.tex is a mechanical, line-for-line copy of the project's single shared preamble as staged in the snapshot (Xy-pic setup plus the native theorem declarations and mathematical macros). Nothing was added, removed or redefined, and the excerpt serves only as preamble, never as a span coordinate.\par\medskip
\noindent\textit{Editorial note.} The source URL is the source-level Stacks Project URL, as required when more than one card is produced from one source; the per-claim tag is carried by claim\_key. The counted statement and proof refer to no labelled item of the chapter and cite no bibliography entry, so this card contributes no external reference stub and needs no bibliography.\par\medskip
\noindent\textit{Allowed standard prerequisites (editorial).}\begin{itemize}
\item Finiteness transfer along a faithfully flat extension: if A to B is faithfully flat, B is Noetherian and B tensor over A for a finite extension A\_1 of A is finite over B, then A\_1 is finite over A. invoked at proof lines 35719-35720 together with the displayed isomorphism, where the author concludes that B over g\textasciicircum{}t B' is a finite B'-module and hence that B' to B is finite. Turns the module-theoretic finiteness of B over g\textasciicircum{}t B' into finiteness of B over B', which is the finiteness assertion of the conclusion.
\item The quotient of a ring by the ideal generated by a nonzerodivisor behaves well under the identification of the two quotients, so that the induced isomorphism of modules may be read modulo f and g. used at proof lines 35714-35718, where the composite A/f\textasciicircum{}tA' -> f\textasciicircum{}tA/f\textasciicircum{}nA' -> g\textasciicircum{}tB/g\textasciicircum{}nB' -> B/g\textasciicircum{}tB' is displayed as an isomorphism of A'/f\textasciicircum{}nA'-modules. Supplies the explicit isomorphism whose compatibility with the ring structures and with the given isomorphism phi' is checked at the end of the proof.
\item Bookkeeping of the ideal inclusions f\textasciicircum{}tA contained in A', f\textasciicircum{}tA' contained in f\textasciicircum{}tA and g\textasciicircum{}tB' contained in N, obtained from A'\_f = A\_f and from the compatibility of phi' with f and g. used at proof lines 35703-35711. Makes B' a subring of B under the definition B = g\textasciicircum{}\{-t\}N and makes the product N N contained in g\textasciicircum{}t N, which is what produces the multiplication on B.
\end{itemize}\medskip
\begin{lemma}
\label{lemma-construct-extension}
Let $A' \subset A$ be an extension of rings. Let $f \in A'$.
Assume that (a) $A$ is finite over $A'$, (b) $f$ is a nonzerodivisor
on $A$, and (c) $A'_f = A_f$. Then there exists an integer $n_0 > 0$
such that for all $n \geq n_0$ the following is true: given a ring
$B'$, a nonzerodivisor $g \in B'$, and an isomorphism
$\varphi' : A'/f^n A' \to B'/g^n B'$ with $\varphi'(f) \equiv g$, there is a
finite extension $B' \subset B$ and an isomorphism
$\varphi : A/fA \to B/gB$ compatible with $\varphi'$.
\end{lemma}

\begin{proof}
Since $A$ is finite over $A'$ and since $A'_f = A_f$ we can
Cchoose $t > 0$ such that $f^t A \subset A'$. Set $n_0 = 2t$.
Given $n, B', g, \varphi'$ as in the statement of the lemma,
denote $N \subset B'$ the set of elements
$b \in B'$ such that $b \bmod g^nB' \in \varphi'(f^tA)$.
Set $B = g^{-t}N$. As $f^tA' \subset f^tA$ and $\varphi'$ sends $f$
to $g$ we have $g^tB' \subset N$, hence $B' \subset B$. Since
$f^tA \cdot f^tA \subset f^t \cdot f^tA$ and $\varphi'$ sends
$f$ to $g$, we see that $N \cdot N \subset g^t N$.
Hence we obtain a multiplication on $B$ extending the
multiplication of $B'$. We have an isomorphism of $A'/f^nA'$-modules
$$
A/f^tA' \xrightarrow{f^t} f^tA/f^nA' \xrightarrow{\varphi'}
g^tB/g^nB' \xrightarrow{g^{-t}} B/g^tB'
$$
where the module structures on the right are defined using $\varphi'$.
Since $A/f^tA'$ is a finite $A'$-module, we conclude that $B/g^tB'$
is a finite $B'$-module and hence we see that $B' \to B$ is finite.
Finally, we leave it to the reader to see that the displayed isomorphism
of modules sends $fA$ into $gB$ and induces an isomorphism of rings
$\varphi : A/fA \to B/gB$ compatible with $\varphi'$ (it even induces an
isomorphism $A/f^tA \to B/g^tB$ but we don't need this).
\end{proof}

\end{document}
